\chapter{Dynamic Games}\label{ch:b2:dynamic-games}

In the \omterm{def:b2:dynamic-games:ultimatum}{ultimatum game} one player proposes how to split a sum of money and the other accepts, in which case
the split is carried out, or rejects, in which case both get nothing. A responder who cares only about
money accepts any positive offer, so a proposer who knows this offers almost nothing. The chart shows what
happened when anthropologists and economists ran the game in eighteen groups from fifteen small-scale
societies, for stakes of about a day or two's wages.

\begin{openingchart}{Ultimatum game in eighteen groups from fifteen small-scale societies: mean offer and
  rejection rate, per cent of the stake. Groups sorted by mean offer; two groups have no rejection
  rate.}{Henrich et al. (2001), Table~1}
\begin{tikzpicture}
\begin{axis}[omrecord, width=0.72\linewidth, height=0.85\linewidth, xmin=0, xmax=65, ymin=0.4,
    ymax=18.6, ymajorgrids=false, clip=false, xlabel={Per cent of the stake},
    ytick=data, yticklabels from table={figdata/bachelor-2/12-dynamic-games/out/ultimatum.csv}{group},
    y tick label style={font=\footnotesize}, legend style={at={(0.5,-0.12)}, anchor=north,
    legend columns=2, /tikz/every even column/.append style={column sep=10pt}}]
  \addplot[xbar, bar width=6pt, fill=omDef!55, draw=omDef] table[x=offer, y=pos]
    {figdata/bachelor-2/12-dynamic-games/out/ultimatum.csv};
  \addplot[only marks, mark=diamond*, mark size=2.4pt, omThm] table[x=rejection, y=pos]
    {figdata/bachelor-2/12-dynamic-games/out/ultimatum.csv};
  \legend{Mean offer, Rejection rate}
\end{axis}
\end{tikzpicture}
\end{openingchart}

% ledger: b2:henrich2001-t1, b2:henrich2001-text
% figdata: ultimatum_summary (lowest 26, highest 58, median 40.5, zero_rejection 3, groups 18)
No group's mean offer is close to zero: the lowest is 26 per cent, the highest 58, and half the groups offer
more than 40 per cent. Rejections are rare in some groups and frequent in others: the Tsimane, Achuar and
Ache rejected none of 70, 16 and 51 offers, while the Gnau rejected 10 of 25.

This chapter studies games in which players move in turn and see what was done before. Their solution,
\omterm{def:b2:dynamic-games:spe}{subgame-perfect equilibrium}, rules out threats that a player would not carry out, and it gives sharp
predictions for bargaining: who gets what, and that agreement is immediate. The \omterm{def:b2:dynamic-games:ultimatum}{ultimatum game} is the
simplest bargaining game, and the chart is where its prediction fails.

\begin{recall}
From \cref{ch:b2:static-games}: \omterm{def:b2:static-games:normal-form}{normal-form games}, \omterm{def:b2:static-games:best-response}{best responses} and \omterm{def:b2:static-games:nash-equilibrium}{Nash equilibrium}. From the Year 1
volume: the leader--follower (Stackelberg) model of quantity competition, with linear demand. From
\cref{ch:b2:intertemporal-choice}: \omterm{def:b2:intertemporal-choice:discount-factor}{discount factors}.
\end{recall}

\begin{notation}[Symbols of this chapter]\label{not:b2:dynamic-games:symbols}
In bargaining, the pie has size 1, $x_i$ is player $i$'s share when she proposes, and $\beta_i$ is her \omterm{def:b2:intertemporal-choice:discount-factor}{discount
factor} per round of bargaining, the value today of one dalen received one round later. In the entry model,
$q_I$ and $q_E$ are the outputs of the incumbent and of the entrant, $F$ is the entrant's fixed cost, and demand
is linear with $a - c = 12$.
\end{notation}

\section{Games in extensive form}\label{sec:b2:dynamic-games:extensive}

\begin{definition}[Extensive-form game]\label{def:b2:dynamic-games:extensive-form}
An \emph{extensive-form game}\index{extensive-form game} is described by a tree: each node where a decision
is taken is labelled with the player who moves there, each branch with an action, and each end node with
the players' payoffs. It has \emph{perfect information} if every player, when she moves, knows every
action taken before.
\end{definition}

A strategy in an \omterm{def:b2:dynamic-games:extensive-form}{extensive-form game} is a complete plan: it says what the player does at \emph{every}
node where she moves, including nodes that will not be reached if she follows it. Take an industry with
one incumbent firm and a potential entrant. The entrant moves first: it stays out, and the incumbent keeps
its monopoly, or enters. After entry the incumbent either accommodates, sharing the market, or fights with
a price war that hurts both.

\begin{omfigure}
\begin{minipage}[b]{0.48\linewidth}\centering
\begin{forest}
for tree={circle, fill=black, inner sep=1.4pt, s sep=18mm, l sep=13mm, edge={thick},
  font=\footnotesize}
[, label={above:Entrant}
  [{$(0, 9)$}, fill=none, rectangle, inner sep=1pt,
    edge label={node[midway, left, xshift=-3pt]{Out}}]
  [, label={right:Incumbent}, edge label={node[midway, right, xshift=3pt]{In}}
    [{$(2, 4)$}, fill=none, rectangle, inner sep=1pt,
      edge label={node[midway, left, xshift=-3pt]{Accommodate}}]
    [{$(-1, 1)$}, fill=none, rectangle, inner sep=1pt,
      edge label={node[midway, right, xshift=3pt]{Fight}}]
  ]
]
\end{forest}
\end{minipage}\hfill
\begin{minipage}[b]{0.48\linewidth}\centering
\begin{tikzpicture}[econ, font=\footnotesize]
  \node[circle, fill=black, inner sep=1.4pt, label={above:Proposer}] (p) at (0,0) {};
  \draw[thick] (p) -- (-1.6,-1.0);
  \draw[thick] (p) -- (1.6,-1.0);
  \draw[thick, densely dotted] (212:1.5) arc[start angle=212, end angle=328, radius=1.5];
  \node[anchor=east] at (-1.6,-1.0) {0};
  \node[anchor=west] at (1.6,-1.0) {1};
  \node[circle, fill=black, inner sep=1.4pt, label={right:Responder}] (r) at (0.4,-1.7) {};
  \draw[thick] (p) -- node[pos=0.35, right] {$x$} (r);
  \draw[thick] (r) -- node[midway, left, xshift=-2pt] {Accept} (-0.5,-2.9);
  \draw[thick] (r) -- node[midway, right, xshift=2pt] {Reject} (1.3,-2.9);
  \node[below] at (-0.5,-2.9) {$(1 - x, x)$};
  \node[below] at (1.3,-2.9) {$(0, 0)$};
\end{tikzpicture}
\end{minipage}

{\small Two games in extensive form. Left: the entry game; payoffs are the entrant's first and the
incumbent's second. Right: the \omterm{def:b2:dynamic-games:ultimatum}{ultimatum game} over a pie of 1; the proposer chooses any offer $x$ between 0
and 1, the arc marking the continuum of offers, and the responder accepts or rejects it.}
\end{omfigure}

The entry game has two Nash equilibria in pure strategies. In one, the entrant enters and the incumbent
accommodates: payoffs $(2, 4)$. In the other, the entrant stays out and the incumbent's strategy is ``fight
if it enters'': payoffs $(0, 9)$. The second is a \omterm{def:b2:static-games:nash-equilibrium}{Nash equilibrium}, since against that plan staying out is
the entrant's \omterm{def:b2:static-games:best-response}{best response}, and the incumbent loses nothing by planning to fight at a node that is never
reached. But the plan is not believable: if entry happened, fighting would earn the incumbent 1 instead of
4.

\section{Backward induction}\label{sec:b2:dynamic-games:backward}

\begin{definition}[Backward induction]\label{def:b2:dynamic-games:backward-induction}
\emph{Backward induction}\index{backward induction} solves a finite game of perfect information from its
end: at each last decision node the mover's best action is chosen; each such node is replaced by the
payoffs it leads to; and the process is repeated until the first node is reached.
\end{definition}

In the entry game, the incumbent accommodates after entry, so entry earns the entrant 2 against 0 for
staying out: it enters. \omterm{def:b2:dynamic-games:backward-induction}{Backward induction} picks the first of the two Nash equilibria and discards the
second.

\begin{definition}[Credible threat]\label{def:b2:dynamic-games:credible-threat}
A threat to take an action at some node is a \emph{credible threat}\index{credible threat} if, were that
node reached, taking the action would be a \omterm{def:b2:static-games:best-response}{best response} for the player who threatens.
\end{definition}

``Fight if you enter'' is not a \omterm{def:b2:dynamic-games:credible-threat}{credible threat}. The idea of ruling out threats that are not credible is
made general by asking for equilibrium play not only in the game as a whole but in every part of it that
could be reached.

\section{Subgame-perfect equilibrium}\label{sec:b2:dynamic-games:spe}

\begin{definition}[Subgame]\label{def:b2:dynamic-games:subgame}
A \emph{subgame}\index{subgame} of a game of perfect information is the part of the tree that starts at one
decision node and contains everything that follows it. The whole game is a subgame of itself.
\end{definition}

\begin{definition}[Subgame-perfect equilibrium]\label{def:b2:dynamic-games:spe}
A \emph{subgame-perfect equilibrium}\index{subgame-perfect equilibrium} is a profile of strategies whose
restriction to every \omterm{def:b2:dynamic-games:subgame}{subgame} is a \omterm{def:b2:static-games:nash-equilibrium}{Nash equilibrium} of that \omterm{def:b2:dynamic-games:subgame}{subgame}.
\end{definition}

\begin{proposition}[Backward induction gives subgame-perfect equilibria]\label{prop:b2:dynamic-games:backward}
In a finite game of perfect information, the strategies found by \omterm{def:b2:dynamic-games:backward-induction}{backward induction} form a
\omterm{def:b2:dynamic-games:spe}{subgame-perfect equilibrium}. If no player has the same payoff at two different end nodes, it is the only
one.
\end{proposition}

\begin{proof}
By induction on the length of the longest path in a \omterm{def:b2:dynamic-games:subgame}{subgame}. A \omterm{def:b2:dynamic-games:subgame}{subgame} with one decision node is a choice by
one player, and the best action is the only \omterm{def:b2:static-games:nash-equilibrium}{Nash equilibrium} (with distinct payoffs). Suppose the claim
holds for all \omterm{def:b2:dynamic-games:subgame}{subgames} shorter than $n$ and take one of length $n$, starting with a move by player $i$. Every
\omterm{def:b2:dynamic-games:subgame}{subgame} that follows her move is shorter, so by the induction hypothesis the backward-induction strategies
are an equilibrium there and fix the payoffs each of her actions leads to. \omterm{def:b2:dynamic-games:backward-induction}{Backward induction} gives her the
action with the highest of these payoffs, so no player can gain by deviating, in the \omterm{def:b2:dynamic-games:subgame}{subgame} or after: the
profile is a \omterm{def:b2:static-games:nash-equilibrium}{Nash equilibrium} of it. With distinct payoffs every choice is strict, so the profile is unique.
\end{proof}

A finite game of perfect information therefore always has a \omterm{def:b2:dynamic-games:spe}{subgame-perfect equilibrium} in pure
strategies, which the method computes. Its prediction can be stark. In the \omterm{def:b2:dynamic-games:ultimatum}{ultimatum game}, the responder
accepts every positive offer, since something beats nothing; with money divisible down to a smallest unit,
the proposer offers that unit or nothing (\cref{exo:b2:dynamic-games:6}). The opening chart is the test.

\section{Commitment}\label{sec:b2:dynamic-games:commitment}

\begin{definition}[Strategic commitment]\label{def:b2:dynamic-games:commitment}
A \emph{strategic commitment}\index{strategic commitment} is an action taken early in a game that cannot be
reversed and that changes the later payoffs of the player who takes it, so that a different action becomes
her \omterm{def:b2:static-games:best-response}{best response} later.
\end{definition}

A commitment works through the other players: by limiting her own future choices, a player changes what
they expect her to do, and so what they do. The leader--follower model of the Year 1 volume is the example:
with demand $p = a - Q$ and $a - c = 12$, a leader that fixes its output first produces 6 and earns 18, more than the 16
each firm earns when both choose at once (4 each). Once the leader's 6 units are produced, the follower can only
make the best of them.

\begin{definition}[Entry deterrence]\label{def:b2:dynamic-games:entry-deterrence}
\emph{Entry deterrence}\index{entry deterrence} is an action by an incumbent firm, taken before a potential
entrant decides, that makes entry unprofitable so that the entrant stays out.
\end{definition}

Let the incumbent choose its output $q_I$ first, by building capacity that it will use. The entrant then
decides whether to enter, at a fixed cost $F$, and if it enters chooses $q_E$. Its \omterm{def:b2:static-games:best-response}{best response} is $q_E = (12 -
q_I)/2$, which earns it $\big((12 - q_I)/2\big)^2 - F$; it stays out when this is not positive.

\begin{omfigure}
\begin{tikzpicture}[econ]
\begin{axis}[omrecord, grid=none, xmin=0, xmax=12.5, ymin=0, ymax=6.8, width=0.62\linewidth,
    height=0.42\linewidth, xlabel={$q_I$: incumbent's output}, ylabel={$q_E$: entrant's output},
    xtick={0,4,6,8,12}, ytick={0,2,3,6}]
  \addplot[omDef, very thick, domain=0:8, samples=2] {(12 - x)/2};
  \addplot[omDef, very thick, domain=8:12, samples=2] {0};
  \addplot[omDef, densely dotted, domain=8:12, samples=2, forget plot] {(12 - x)/2};
  \draw[guide] (axis cs:8,0) -- (axis cs:8,2);
  \node[eqpoint] at (axis cs:6,3) {};
  \node[eqpoint] at (axis cs:8,0) {};
  \node[anchor=south west, font=\footnotesize] at (axis cs:6.1,3.05) {accommodate: leader at 6};
  \node[anchor=south west, font=\footnotesize] at (axis cs:8.1,0.15) {deter: 8};
  \node[anchor=north west, font=\footnotesize, omDef] at (axis cs:0.4,3.2)
    {entrant's best response};
\end{axis}
\end{tikzpicture}

{\small The entrant's \omterm{def:b2:static-games:best-response}{best response} to the incumbent's committed output, with a fixed cost $F = 4$. Below
$q_I = 8$ the entrant enters and produces $(12 - q_I)/2$; from 8 on its profit would not cover the fixed cost
and it stays out. The incumbent compares the leader's point, where it earns 18, with deterrence at 8, where it
earns 32.}
\end{omfigure}

With $F = 4$ the entrant stays out once $q_I \ge 8$. The incumbent can accommodate, choosing the leader's output 6, at
which the entrant enters with 3 and the price is 3 above cost: profit 18. Or it can deter, choosing 8, at which the
price is 4 above cost and the entrant stays out: profit 32. It deters. With a larger fixed cost the monopoly
output 6 already deters, and entry is blockaded; with a much smaller one, deterrence needs so much output that
accommodating pays more (\cref{exo:b2:dynamic-games:11}). The output deters only because it is committed: an
incumbent that could cut back after entry would do so, and the entrant, knowing it, would enter.

\section{Bargaining by alternating offers}\label{sec:b2:dynamic-games:bargaining}

\begin{definition}[Ultimatum game]\label{def:b2:dynamic-games:ultimatum}
The \emph{ultimatum game}\index{ultimatum game} is a two-player game in which a proposer offers a division
of a fixed sum and a responder either accepts, and the division is carried out, or rejects, and both
receive nothing.
\end{definition}

\begin{definition}[Alternating-offers bargaining model]\label{def:b2:dynamic-games:alternating-offers}
In the \emph{alternating-offers bargaining model}\index{alternating-offers bargaining model} two players
divide a pie of size 1. In round 1 player 1 proposes a division and player 2 accepts or rejects it; after a
rejection, player 2 proposes in round 2 and player 1 responds, and so on. A share $s$ agreed in round $t$ is
worth $\beta_i^{\,t-1} s$ to player $i$, with $0 < \beta_i < 1$. Bargaining may end after a fixed number of rounds, with
nothing for either player, or go on without limit.
\end{definition}

The \omterm{def:b2:dynamic-games:ultimatum}{ultimatum game} is the model with one round. With two rounds, player 2 will take the whole pie in round 2,
which is worth $\beta_2$ to her today; so in round 1 she accepts any offer giving her $\beta_2$, and player 1 keeps $1 -
\beta_2$. Each extra round adds one step of the same argument. The listing runs the steps for \omterm{def:b2:intertemporal-choice:discount-factor}{discount factors}
$\beta_1 = 0.95$ and $\beta_2 = 0.90$.

\begin{omfigure}
\begin{forest}
for tree={circle, fill=black, inner sep=1.4pt, s sep=22mm, l sep=11mm, edge={thick}, font=\footnotesize}
[, label={above:Player 1 proposes}
  [, label={right:2}, edge label={node[midway, right, xshift=2pt]{$(x, 1 - x)$}}
    [{$(x, 1 - x)$}, fill=none, rectangle, inner sep=1pt,
      edge label={node[midway, left, xshift=-3pt]{Accept}}]
    [, label={right:2 proposes}, edge label={node[midway, right, xshift=3pt]{Reject}}
      [, label={right:1}, edge label={node[midway, right, xshift=2pt]{$(1 - y, y)$}}
        [{$(\beta_1(1 - y), \beta_2 y)$}, fill=none, rectangle, inner sep=1pt,
          edge label={node[midway, left, xshift=-3pt]{Accept}}]
        [{$(0, 0)$}, fill=none, rectangle, inner sep=1pt,
          edge label={node[midway, right, xshift=3pt]{Reject}}]
      ]
    ]
  ]
]
\end{forest}

{\small Alternating offers with two rounds. Player 1 proposes a division, which player 2 accepts or rejects;
after a rejection player 2 proposes and player 1 responds. Payoffs are valued in round 1: a share received in round 2 is
discounted by each player's $\beta_i$. \omterm{def:b2:dynamic-games:backward-induction}{Backward induction}: in round 2 player 2 offers $y = 1$, which player 1 accepts
(or is indifferent about); so in round 1 player 2 accepts any $1 - x \ge \beta_2$, and player 1 offers $x = 1 - \beta_2$.}
\end{omfigure}

\omcode{code/bachelor-2/12-dynamic-games/listings/bargaining.py}{1}{20}{Alternating offers with a final round,
solved backwards, for longer and longer deadlines.}\label{lst:b2:dynamic-games:bargaining}

\omoutput{code/bachelor-2/12-dynamic-games/outputs/bargaining.txt}{With a short deadline the last proposer gets
most of the pie; as the deadline recedes, the split settles down.}\label{out:b2:dynamic-games:bargaining}

% code: listings/bargaining.py
With a deadline of one or three rounds, player 1 makes the last offer and gets most of the pie; with two or
four, player 2 does. The deadline decides, which is unsatisfactory when real negotiations have none. As the
number of rounds grows, the effect of the last round is discounted away and the share settles at 0.6897.
That limit is the equilibrium of the game without a deadline.

\begin{theorem}[Alternating offers without a deadline]\label{thm:b2:dynamic-games:rubinstein}
The \omterm{def:b2:dynamic-games:alternating-offers}{alternating-offers bargaining model} without a deadline has a unique \omterm{def:b2:dynamic-games:spe}{subgame-perfect equilibrium}. Each
player, when she proposes, offers the other exactly what the other would get by rejecting and proposing
next round; each accepts any offer at least that large. Player 1's offer is accepted at once and gives her
\[
x_1 = \frac{1 - \beta_2}{1 - \beta_1 \beta_2}.
\]
\end{theorem}

\begin{proof}[Proof that these strategies form an equilibrium, and the formula]
Suppose each player, when she proposes, asks for a share $x_i$ that does not depend on the history, and accepts
any offer worth at least what she would get by rejecting. Rejecting an offer from player 1 earns player 2 her
own share next round, worth $\beta_2 x_2$ today, so player 1 must leave her that: $1 - x_1 = \beta_2 x_2$. In the same
way $1 - x_2 = \beta_1 x_1$. Substituting, $x_1 = 1 - \beta_2(1 - \beta_1 x_1)$, which gives the formula. No player gains
from a single deviation at any node: asking for more is rejected and leaves her, as next round's responder,
$\beta_i(1 - x_j) = \beta_i^2 x_i < x_i$ in today's value; asking for less gives away share for nothing; rejecting an
offer that meets her threshold earns her no more. Because payoffs are discounted, a strategy that gains from no
single deviation gains from none. Uniqueness is \admitted; \cref{exo:b2:dynamic-games:10} sketches it.
\end{proof}

The theorem's predictions are sharp. Agreement is immediate, so no pie is wasted. Patience is power: with
$\beta_1 = 0.95$ and $\beta_2 = 0.90$, player 1 gets 0.69. With equal patience the first proposer gets $1/(1 + \beta)$, a
small advantage that vanishes as the time between offers shrinks and $\beta$ approaches 1, leaving an even
split (\cref{exo:b2:dynamic-games:3}).

\begin{model}{Alternating-offers bargaining}
\assumes Two players who maximise discounted shares, each knowing the other's \omterm{def:b2:intertemporal-choice:discount-factor}{discount factor}, with no
outside options.
\predicts Immediate agreement; a larger share for the more patient player; a first-mover advantage that
vanishes as offers come faster.
\tested Most directly through its one-round case, the \omterm{def:b2:dynamic-games:ultimatum}{ultimatum game} (Henrich et al. 2001; Andersen et al.
2011).
\breaks The model cannot produce a delay: every negotiation that takes time, or breaks down for a while,
lies outside it. Delays appear in versions in which each side is unsure of the other's patience or costs.
\end{model}

\section{Ultimatum experiments: where the model breaks}\label{sec:b2:dynamic-games:ultimatum}

The \omterm{def:b2:dynamic-games:ultimatum}{ultimatum game} has been run with student subjects in many countries, and its prediction --- a tiny
offer, accepted --- fails there too.
% ledger: b2:henrich2001-text
Offers in industrial societies typically average close to 44 per cent of the stake with a mode at 50, and
offers below 20 per cent are rejected with probability 0.40 to 0.60 (Henrich et al. 2001, summarising earlier
studies). A responder who rejects 20 per cent gives up money to deny the proposer 80; a proposer who expects
this offers more than the minimum, and may also care about fairness himself.

\begin{record}[Ultimatum offers in fifteen societies]
% ledger: b2:henrich2001-t1, b2:henrich2001-text
Henrich et al. (2001) ran the game in eighteen groups from fifteen small-scale societies, mostly with money,
for stakes worth about a day or two's wages. Mean offers ranged from 26 per cent (Machiguenga, Peru) to 58 per
cent (Lamelara, Indonesia), wider than among students. Rejections varied as much: the Tsimane of Bolivia
rejected none of 70 offers, although nearly half of the offers were 30 per cent or less, while the Au and Gnau
of Papua New Guinea rejected 8 of 30 and 10 of 25. No group behaved as the \omterm{def:b2:dynamic-games:spe}{subgame-perfect equilibrium} with
money-maximising players predicts.
\end{record}

\begin{model}{Subgame-perfect equilibrium with money-maximising players}
\assumes Each player maximises her own money payoff, and this is common knowledge.
\predicts In the \omterm{def:b2:dynamic-games:ultimatum}{ultimatum game}, the smallest possible offer, accepted.
\tested Ultimatum experiments with students and in small-scale societies (Henrich et al. 2001); at stakes up to
1{,}600 hours of work (Andersen et al. 2011).
\breaks Mean offers of 26 to 58 per cent; rejections of positive offers, which fall significantly only at the highest
stakes.
\end{model}

The failure is of the payoff assumption, not of \omterm{def:b2:dynamic-games:backward-induction}{backward induction} itself: if a responder dislikes unequal
splits enough to reject small offers, a proposer reasoning backwards should offer more, as the data show. The
open question is how large the stakes must be before money dominates.

\begin{debate}[Do higher stakes change ultimatum behaviour?]
% ledger: b2:aeghl2011-design, b2:aeghl2011-t2
Earlier experiments found that responders did not markedly change their behaviour as the stakes rose
(Slonim and Roth 1998; Cameron 1999; both as summarised by Andersen et al. 2011). Andersen et al. (2011) ran the game in
eight villages of Meghalaya, in Northeast India, for 20, 200, 2{,}000 or 20{,}000 rupees, from 1.6 to 1{,}600 hours of
work, with instructions that drew many low offers. Rejection rates were 36 per cent at 20 rupees, 43 at 200, 28 at
2{,}000 and 4 at 20{,}000; only the drop to the highest stakes was statistically significant. One reading: fairness matters
at the stakes of most experiments and of much everyday bargaining, and money dominates only far beyond them.
Another: the highest stakes were played by only 24 responders, in one region, and the earlier results stand for the
stakes people usually face.
\end{debate}

\section{Exercises}\label{sec:b2:dynamic-games:exercises}

\begin{exercise}[$\star$]\label{exo:b2:dynamic-games:1}
Player 1 chooses Out, ending the game with payoffs $(2, 2)$, or In. After In, player 2 chooses Low, ending with
$(1, 3)$, or High, after which player 1 chooses $a$, giving $(4, 1)$, or $b$, giving $(0, 5)$. Payoffs are player 1's first.
Solve the game by \omterm{def:b2:dynamic-games:backward-induction}{backward induction}.
\end{exercise}

\begin{exercise}[$\star$]\label{exo:b2:dynamic-games:2}
Read the opening chart. What are the lowest and highest mean offers? How many groups made mean offers below
30 per cent, and how many rejected no offer at all?
\end{exercise}

\begin{exercise}[$\star$]\label{exo:b2:dynamic-games:3}
Predict, then compute: two equally patient players, $\beta_1 = \beta_2 = 0.9$, bargain by alternating offers without a
deadline. What does the first proposer get? What happens to her share as $\beta$ approaches 1?
\end{exercise}

\begin{exercise}[$\star$]\label{exo:b2:dynamic-games:4}
A union and an employer bargain for six weeks before agreeing. Does this refute the alternating-offers model?
What would a test of the model need to observe?
\end{exercise}

\begin{exercise}[$\star\star$]\label{exo:b2:dynamic-games:5}
Write the entry game in normal form, find all its Nash equilibria in pure strategies, and say which are
\omterm{def:b2:dynamic-games:subgame}{subgame} perfect.
\end{exercise}

\begin{exercise}[$\star\star$]\label{exo:b2:dynamic-games:6}
In an \omterm{def:b2:dynamic-games:ultimatum}{ultimatum game} over 10 coins, offers must be whole coins. Find every \omterm{def:b2:dynamic-games:spe}{subgame-perfect equilibrium}.
\end{exercise}

\begin{exercise}[$\star\star$]\label{exo:b2:dynamic-games:7}
Solve the alternating-offers game with three rounds and $\beta_1 = \beta_2 = 0.9$ by \omterm{def:b2:dynamic-games:backward-induction}{backward induction}. Compare with
\cref{exo:b2:dynamic-games:3}.
\end{exercise}

\begin{exercise}[$\star\star$]\label{exo:b2:dynamic-games:8}
Use the rejection rates of Andersen et al.\ (2011) in the debate box. What does the \omterm{def:b2:dynamic-games:spe}{subgame-perfect equilibrium}
with money-maximising responders predict for them? Do the data move toward the prediction as the stakes rise,
and steadily?
\end{exercise}

\begin{exercise}[$\star\star$]\label{exo:b2:dynamic-games:9}
In the entry model with $F = 4$, check that output 8 deters entry and that the incumbent prefers it to the
leader's output. Why must the output be chosen before the entrant decides and be impossible to reduce?
\end{exercise}

\begin{exercise}[$\star\star\star$]\label{exo:b2:dynamic-games:10}
Let $M_1$ and $m_1$ be the largest and smallest payoffs player 1 gets in any \omterm{def:b2:dynamic-games:spe}{subgame-perfect equilibrium} of the
alternating-offers game without a deadline, in a \omterm{def:b2:dynamic-games:subgame}{subgame} where she proposes, and $M_2$, $m_2$ the same for player 2.
Show that $m_1 \ge 1 - \beta_2 M_2$ and $M_2 \le 1 - \beta_1 m_1$, and the same with the roles exchanged, and deduce that $M_1 = m_1
= x_1$.
\end{exercise}

\begin{exercise}[$\star\star\star$]\label{exo:b2:dynamic-games:11}
In the entry model, find the fixed cost $F$ below which the incumbent prefers to accommodate, and the fixed cost
above which entry is blockaded.
\end{exercise}

\begin{exercise}[$\star\star\star$]\label{exo:b2:dynamic-games:12}
An incumbent has a store in each of 20 towns, and in each town in turn a different entrant plays the entry game
of this chapter. Use \omterm{def:b2:dynamic-games:backward-induction}{backward induction} to show that every entrant enters and the incumbent always
accommodates. Why might a reputation for fighting still pay if the entrants were unsure of the incumbent's
payoffs?
\end{exercise}

\section{Problem: harbour wages in Brenmoor}\label{sec:b2:dynamic-games:problem}

\begin{problem}[{Weekend problem --- the dockworkers' union and the harbour employer of Brenmoor, the capital of
Brenland, an invented economy, bargain over a wage premium: how is it split, why is there no strike, and what
moves the split}]
\label{pb:b2:dynamic-games:1}
The premium to be divided is 1{,}000 dalens per worker per month above what the workers would earn elsewhere.
Each day without agreement the harbour is shut. The employer values a share received one day later at 0.98 of its
value today, the union at 0.95. The employer proposes on day 1, the union on day 2, and so on.

\medskip
\noindent\textbf{Part I --- The cost of a strike.}
\begin{enumerate}
  \item What is a premium of 1{,}000 dalens received tomorrow worth to the union today?
  \item Which side is more patient?
  \item What does one day of disagreement cost each side on a premium worth 1{,}000 to it?
  \item If the union could get the whole premium tomorrow, what is the least it would accept today?
  \item Why does a day of strike cost the employer anything?
  \item Name one thing that could raise the union's \omterm{def:b2:intertemporal-choice:discount-factor}{discount factor}.
\end{enumerate}

\medskip
\noindent\textbf{Part II --- A deadline.} Suppose talks must end on day 2, after the union's offer.
\begin{enumerate}[resume]
  \item What does the union propose on day 2?
  \item What is the least the union accepts on day 1?
  \item What does the employer offer on day 1, and what is the split?
  \item Is there a strike?
  \item Find the split if instead the talks end on day 3, after a second offer by the employer.
  \item What do these two answers show about deadlines?
\end{enumerate}

\medskip
\noindent\textbf{Part III --- No deadline.}
\begin{enumerate}[resume]
  \item Write the two conditions that fix the shares each side asks for when it proposes.
  \item Solve them: what does each side get, in dalens?
  \item Is there a strike in equilibrium?
  \item Find the split if the union proposed first. How large is the first-mover advantage?
  \item A strike fund raises the union's \omterm{def:b2:intertemporal-choice:discount-factor}{discount factor} to 0.97. Find the new split.
  \item The union's members have run down their savings, and its \omterm{def:b2:intertemporal-choice:discount-factor}{discount factor} falls to 0.90. Find the split.
\end{enumerate}

\medskip
\noindent\textbf{Part IV --- An outside option.} After rejecting an offer of the employer, the workers can leave for
other jobs that are worth $O$ dalens of the premium to them; then talks end and the employer gets nothing.
\begin{enumerate}[resume]
  \item With $O = 200$, does the split of Part III change?
  \item With $O = 350$, find the split.
  \item Why does an outside option matter only when it is large enough?
  \item Which gives the union more: the strike fund of question 17 or an outside option of 350?
  \item How many dalens does the strike fund move to the union?
  \item State the result: the split, and how the union's cost of a strike day moves it.
\end{enumerate}
\end{problem}
